\documentclass[12pt,a4paper]{article}

% ── Encodage & langue ───────────────────────────────────────────────────────
\usepackage[T1]{fontenc}
\usepackage[utf8]{inputenc}
\usepackage[french]{babel}

% ── Mise en page ────────────────────────────────────────────────────────────
\usepackage[top=2.5cm, bottom=2.5cm, left=2.5cm, right=2.5cm]{geometry}
\usepackage{parskip}
\setlength{\parindent}{0pt}

% ── Typographie ─────────────────────────────────────────────────────────────
\usepackage{lmodern}
\usepackage{microtype}
\setlength{\emergencystretch}{50pt}  % absorbe les dépassements de \texttt{}

% ── Couleurs & liens ────────────────────────────────────────────────────────
\usepackage[dvipsnames,svgnames]{xcolor}
\definecolor{primary}{HTML}{0f766e}
\definecolor{codebg}{HTML}{f6f8fa}
\definecolor{codefg}{HTML}{24292e}
\definecolor{linkcol}{HTML}{0f766e}

\usepackage[
  colorlinks=true,
  linkcolor=linkcol,
  urlcolor=linkcol,
  citecolor=linkcol,
  bookmarks=true,
  pdfauthor={Finance Tracker},
  pdftitle={Architecture du frontend Angular -- Finance Tracker}
]{hyperref}

% ── Code ────────────────────────────────────────────────────────────────────
\usepackage{listings}
\lstdefinelanguage{TypeScript}{
  keywords={import,export,from,class,interface,type,const,let,var,function,return,
            if,else,for,while,switch,case,break,continue,new,this,super,extends,
            implements,async,await,true,false,null,undefined,void,number,string,
            boolean,any,never,unknown,readonly,private,public,protected,static,
            abstract,enum,namespace,module,declare,as,in,of,default,throw,try,
            catch,finally,delete,typeof,instanceof,keyof,Partial,Pick,Omit,
            Observable,Promise,inject,Component,Injectable,OnInit,Input,Output},
  sensitive=true,
  comment=[l]{//},
  morecomment=[s]{/*}{*/},
  string=[b]",
  morestring=[b]',
  morestring=[b]`,
}
\lstset{
  language=TypeScript,
  backgroundcolor=\color{codebg},
  basicstyle=\ttfamily\footnotesize,
  keywordstyle=\color{primary}\bfseries,
  commentstyle=\color{gray}\itshape,
  stringstyle=\color{BrickRed},
  numbers=left,
  numberstyle=\tiny\color{gray},
  numbersep=8pt,
  tabsize=2,
  showstringspaces=false,
  breaklines=true,
  frame=single,
  rulecolor=\color{gray!30},
  xleftmargin=1.5em,
  framexleftmargin=1.5em,
  captionpos=b,
}
\lstdefinestyle{shell}{
  language=bash,
  backgroundcolor=\color{black!88},
  basicstyle=\ttfamily\small\color{white},
  keywordstyle=\color{white},
  commentstyle=\color{gray!60},
  stringstyle=\color{white},
  numbers=none,
  frame=none,
  xleftmargin=0pt,
  framexleftmargin=0pt,
}

% ── Tableaux ────────────────────────────────────────────────────────────────
\usepackage{booktabs}
\usepackage{tabularx}
\usepackage{longtable}
\usepackage{array}
\newcolumntype{L}[1]{>{\raggedright\arraybackslash}p{#1}}

% ── Titres ──────────────────────────────────────────────────────────────────
\usepackage{titlesec}
\titleformat{\section}
  {\Large\bfseries\color{primary}}
  {\thesection}{1em}{}[\color{primary}\titlerule]
\titleformat{\subsection}
  {\large\bfseries\color{codefg}}
  {\thesubsection}{1em}{}
\titleformat{\subsubsection}
  {\normalsize\bfseries\color{codefg}}
  {\thesubsubsection}{1em}{}

% ── Divers ──────────────────────────────────────────────────────────────────
\usepackage{graphicx}
\usepackage{enumitem}
\usepackage{fancybox}
\usepackage{mdframed}

\mdfdefinestyle{infobox}{
  backgroundcolor=primary!8,
  linecolor=primary,
  linewidth=1pt,
  leftmargin=0pt, rightmargin=0pt,
  innerleftmargin=10pt, innerrightmargin=10pt,
  innertopmargin=8pt, innerbottommargin=8pt,
}

\usepackage{fancyhdr}
\pagestyle{fancy}
\fancyhf{}
\rhead{\textcolor{gray}{\small Finance Tracker -- Architecture Frontend}}
\lhead{\textcolor{gray}{\small Angular 21 + Supabase + Spring Boot}}
\rfoot{\thepage}
\renewcommand{\headrulewidth}{0.4pt}

% ── Page de garde ───────────────────────────────────────────────────────────
\title{
  \vspace{-1cm}
  {\color{primary}\rule{\linewidth}{2pt}}\\[0.5cm]
  {\LARGE\bfseries Finance Tracker}\\[0.3cm]
  {\large\bfseries Architecture du Frontend Angular}\\[0.2cm]
  {\normalsize Documentation technique complète}\\[0.5cm]
  {\color{primary}\rule{\linewidth}{2pt}}
}
\author{}
\date{\today}

% ════════════════════════════════════════════════════════════════════════════
\begin{document}
\maketitle
\thispagestyle{empty}
\tableofcontents
\newpage

% ════════════════════════════════════════════════════════════════════════════
\section{Vue d'ensemble}

Finance Tracker est une application de gestion de finances personnelles.
Le \textbf{frontend} est une Single Page Application (SPA) construite avec
\textbf{Angular 21}, qui consomme une API REST exposée par un backend
\textbf{Spring Boot} et délègue l'authentification à \textbf{Supabase Auth}.

\subsection{Stack technique}

\begin{center}
\begin{tabular}{L{4cm} L{9cm}}
\toprule
\textbf{Technologie} & \textbf{Rôle} \\
\midrule
Angular 21 & Framework SPA, standalone components, lazy loading \\
TypeScript strict & Typage statique complet, aucun \texttt{any} autorisé \\
Angular Router & Navigation côté client, guards, lazy routes \\
Angular HttpClient & Requêtes HTTP, intercepteurs \\
Angular Reactive Forms & Formulaires typés avec validation \\
RxJS & Programmation réactive (Observables) \\
Supabase JS SDK & Authentification côté client \\
Tailwind CSS 4 & Utilitaires CSS \\
chart.js + ng2-charts & Graphique linéaire du tableau de bord \\
SweetAlert2 & Boîtes de dialogue et alertes \\
dayjs & Manipulation des dates \\
\bottomrule
\end{tabular}
\end{center}

\subsection{Périmètre fonctionnel}

L'application couvre sept domaines métier :

\begin{itemize}[itemsep=2pt]
  \item \textbf{Transactions} -- saisie, filtrage, recherche, pagination
  \item \textbf{Catégories} -- gestion des catégories revenus/dépenses
  \item \textbf{Budgets} -- plafonds mensuels par catégorie avec alertes
  \item \textbf{Objectifs} -- épargne progressive avec historique des contributions
  \item \textbf{Dettes} -- suivi du remboursement avec historique des paiements
  \item \textbf{Transactions planifiées} -- récurrences automatisées
  \item \textbf{Calendrier} -- vue mensuelle avec projection des récurrences
  \item \textbf{Import / Export} -- CSV bidirectionnel
  \item \textbf{Tableau de bord} -- agrégats et visualisations
\end{itemize}

% ════════════════════════════════════════════════════════════════════════════
\section{Structure des répertoires}

Le projet suit une architecture en quatre couches :
\texttt{core}, \texttt{shared}, \texttt{layout}, \texttt{features}.

\begin{lstlisting}[language=bash, numbers=none, frame=none, backgroundcolor=\color{codebg}, basicstyle=\ttfamily\small\color{codefg}]
src/app/
  app.ts                       # Composant racine
  app.config.ts                # Configuration applicative (providers)
  app.routes.ts                # Définition des routes

  core/                        # Infrastructure transverse
    auth/
      supabase.client.ts       # Client Supabase (singleton)
    guards/
      auth-guard.ts            # Protège les routes privées
      guest-guard.ts           # Bloque un user connecté sur /login
    interceptors/
      auth.interceptor.ts      # Injecte le Bearer token
    models/
      common.model.ts          # Types partagés (Page<T>, ApiError, enums)
    services/
      auth.service.ts          # Gestion de la session Supabase
      account.service.ts       # Déconnexion globale (logout-all)
    validators/
      custom-validators.ts     # passwordMatchValidator (factory)

  shared/                      # Composants UI réutilisables
    components/
      pagination/              # PaginationComponent (Spring Page)

  layout/                      # Conteneurs visuels
    main-layout/               # Sidebar + router-outlet

  features/                    # Domaines métier
    landing/                   # Page d'accueil publique
    auth/                      # login, signup, forgot-password, reset-password
    dashboard/                 # Tableau de bord
    transactions/              # CRUD transactions
    categories/                # CRUD catégories
    budgets/                   # CRUD budgets
    goals/                     # Objectifs + contributions
    debts/                     # Dettes + paiements
    planned-transactions/      # Récurrences + runs
    calendar/                  # Calendrier mensuel
    import-export/             # Import/export CSV
    profile/                   # Profil utilisateur
    sidebar/                   # Composant navigation (AppSidebar)
\end{lstlisting}

% ════════════════════════════════════════════════════════════════════════════
\section{Configuration applicative (\texttt{app.config.ts})}

Angular 21 n'utilise plus de \texttt{NgModule} racine.
La configuration est déclarative dans \texttt{app.config.ts} :

\begin{lstlisting}
export const appConfig: ApplicationConfig = {
  providers: [
    provideZoneChangeDetection({ eventCoalescing: true }),
    provideRouter(routes),
    provideHttpClient(withInterceptors([authInterceptor])),
    provideCharts(withDefaultRegisterables()),
  ]
};
\end{lstlisting}

\begin{itemize}
  \item \texttt{provideZoneChangeDetection} -- active la détection de changements avec
        coalescing (une seule passe par macro-tâche au lieu d'une par événement).
  \item \texttt{provideRouter(routes)} -- enregistre les routes avec lazy loading.
  \item \texttt{provideHttpClient(withInterceptors([...]))} -- configure
        \texttt{HttpClient} en mode standalone et branche \texttt{authInterceptor}
        sur \textbf{toutes} les requêtes HTTP sortantes.
  \item \texttt{provideCharts} -- enregistre les composants chart.js nécessaires
        à \texttt{ng2-charts}.
  \item Tous les services métier utilisent \texttt{providedIn: 'root'} : ils sont
        des \textbf{singletons} créés à la demande par l'injecteur Angular.
\end{itemize}

% ════════════════════════════════════════════════════════════════════════════
\section{Routage (\texttt{app.routes.ts})}

\subsection{Routes publiques}

\begin{center}
\begin{tabular}{L{5cm} L{6cm} L{3cm}}
\toprule
\textbf{URL} & \textbf{Composant} & \textbf{Guard} \\
\midrule
\texttt{/} & \texttt{LandingPage} & -- \\
\texttt{/login} & \texttt{LoginPage} & \texttt{guestGuard} \\
\texttt{/signup} & \texttt{SignupPage} & \texttt{guestGuard} \\
\texttt{/forgot-password} & \texttt{ForgotPasswordPage} & \texttt{guestGuard} \\
\texttt{/reset-password} & \texttt{ResetPasswordPage} & -- \\
\bottomrule
\end{tabular}
\end{center}

\subsection{Routes privées}

Toutes sous \texttt{/app}, encapsulées dans \texttt{MainLayout} et
protégées par \texttt{authGuard} :

\begin{center}
\begin{tabular}{L{6cm} L{7.5cm}}
\toprule
\textbf{URL} & \textbf{Composant} \\
\midrule
\texttt{/app/dashboard} & \texttt{DashboardPage} \\
\texttt{/app/transactions} & \texttt{TransactionsPage} \\
\texttt{/app/categories} & \texttt{CategoriesPage} \\
\texttt{/app/budgets} & \texttt{BudgetsPage} \\
\texttt{/app/goals} & \texttt{GoalsPage} \\
\texttt{/app/debts} & \texttt{DebtsPage} \\
\texttt{/app/planned} & \texttt{PlannedPage} \\
\texttt{/app/calendar} & \texttt{CalendarPage} \\
\texttt{/app/import-export} & \texttt{ImportExportPage} \\
\texttt{/app/profile} & \texttt{ProfilePage} \\
\bottomrule
\end{tabular}
\end{center}

Chaque route utilise \texttt{loadComponent()} : le bundle JavaScript du composant
est téléchargé \textbf{uniquement} lorsque l'utilisateur accède à la page
(\textit{lazy loading}).

\begin{lstlisting}
{
  path: 'transactions',
  loadComponent: () =>
    import('./features/transactions/pages/transactions-page/transactions-page')
      .then(m => m.TransactionsPage)
}
\end{lstlisting}

% ════════════════════════════════════════════════════════════════════════════
\section{Couche Core -- Infrastructure}

\subsection{Client Supabase (\texttt{supabase.client.ts})}

Le client Supabase est instancié \textbf{une seule fois} au démarrage de l'app,
avec la clé publique \texttt{anon} (jamais la \texttt{service\_role}) :

\begin{lstlisting}
export const supabase = createClient(
  environment.supabaseUrl,
  environment.supabaseAnonKey,
  {
    auth: {
      lock: (_name, _timeout, fn) => fn(),
    },
  }
);
\end{lstlisting}

L'option \texttt{lock} désactive le verrou interne de Supabase, qui utilise
normalement la Web Locks API (\texttt{navigator.locks}). Cela permet à
\texttt{AuthService} de lire la session de façon \textbf{synchrone} depuis
le cache, sans attendre une Promise à chaque requête HTTP.

\subsection{Service d'authentification (\texttt{auth.service.ts})}

\begin{lstlisting}
@Injectable({ providedIn: 'root' })
export class AuthService {
  private session$ = new BehaviorSubject<Session | null>(null);
  private initialized = false;

  constructor(private zone: NgZone) {
    supabase.auth.getSession().then(({ data }) => {
      this.session$.next(data.session);
      this.initialized = true;
    });
    supabase.auth.onAuthStateChange((_event, session) => {
      this.zone.run(() => this.session$.next(session));
    });
  }

  getAccessToken(): string | null {
    return this.session$.value?.access_token ?? null;
  }

  isAuthenticated(): boolean {
    return this.session$.value !== null;
  }
}
\end{lstlisting}

\textbf{Mécanismes clés :}
\begin{itemize}
  \item \texttt{BehaviorSubject<Session | null>} -- émet la session courante à chaque
        changement ; valeur initiale \texttt{null} (non connecté).
  \item \texttt{waitForInit()} -- les guards attendent que \texttt{getSession()} ait
        répondu avant de décider si l'utilisateur est authentifié.
  \item \texttt{getAccessToken()} -- lecture \textbf{synchrone} du token depuis
        \texttt{BehaviorSubject.value}. Supabase met à jour le cache localement ;
        l'intercepteur n'a pas besoin de souscrire à un Observable.
  \item \texttt{NgZone.run()} -- les événements Supabase arrivent hors de la zone
        Angular ; cette ligne force la détection de changements lors du changement
        de session.
\end{itemize}

\subsection{Garde d'authentification (\texttt{auth-guard.ts})}

\begin{lstlisting}
export const authGuard: CanActivateFn = async () => {
  const auth   = inject(AuthService);
  const router = inject(Router);
  await auth.waitForInit();
  if (!auth.isAuthenticated())
    return router.createUrlTree(['/login']);
  return true;
};
\end{lstlisting}

Appliqué au groupe de routes \texttt{/app/**}. Si la session n'est pas encore chargée
au démarrage (premier chargement de l'app), le guard \textbf{attend} la résolution
de \texttt{waitForInit()} avant de statuer.

\subsection{Garde invité (\texttt{guest-guard.ts})}

Empêche un utilisateur \textbf{déjà connecté} d'accéder à \texttt{/login} ou
\texttt{/signup}. Il est redirigé vers \texttt{/app/dashboard}.

\begin{lstlisting}
export const guestGuard: CanActivateFn = async () => {
  const auth   = inject(AuthService);
  const router = inject(Router);
  await auth.waitForInit();
  if (auth.isAuthenticated())
    return router.createUrlTree(['/app/dashboard']);
  return true;
};
\end{lstlisting}

\subsection{Intercepteur HTTP (\texttt{auth.interceptor.ts})}

\begin{lstlisting}
export const authInterceptor: HttpInterceptorFn = (req, next) => {
  const token = inject(AuthService).getAccessToken();
  if (!token) return next(req);
  return next(req.clone({
    setHeaders: { Authorization: `Bearer ${token}` }
  }));
};
\end{lstlisting}

\textbf{Fonctionnement :}
\begin{enumerate}
  \item Chaque requête HTTP passe par cet intercepteur (enregistré globalement).
  \item Le token JWT est lu \textbf{synchronement} depuis le cache d'AuthService.
  \item Si un token existe, la requête est clonée (immuabilité) avec l'en-tête
        \texttt{Authorization: Bearer <token>}.
  \item Si aucun token n'existe (utilisateur non connecté), la requête passe telle quelle.
  \item Le backend Spring Boot valide ce JWT sur toutes les routes protégées.
\end{enumerate}

\subsection{Architecture du backend vue depuis le frontend}

Même si ce dépôt contient uniquement le frontend Angular, le contrat HTTP et les
types manipulés permettent de reconstituer précisément la \textbf{chaîne de traitement}
côté backend. Le modèle le plus cohérent avec les endpoints utilisés est une
architecture Spring Boot classique en couches :

\begin{center}
\begin{tabularx}{\linewidth}{L{3.2cm} X}
\toprule
\textbf{Couche backend} & \textbf{Responsabilité} \\
\midrule
Security filter chain & Intercepte la requête entrante, lit l'en-tête
\texttt{Authorization}, valide le JWT Supabase, refuse la requête si le token
est absent ou invalide sur une route protégée. \\
\midrule
Controller REST & Expose les routes \texttt{/v1/**}, reçoit le JSON, valide les
données entrantes, convertit la requête HTTP en appel métier. \\
\midrule
Service métier & Porte les règles métier : création, recherche, contrôle des
droits, calculs, déclenchement d'actions liées, orchestration entre plusieurs
repositories si nécessaire. \\
\midrule
Repository / persistance & Charge et sauvegarde les entités en base, applique les
filtres, tris, paginations et recherches demandés par le service. \\
\midrule
DTO / sérialisation & Transforme les objets métier en JSON de réponse stable pour
le frontend : listes, pages, messages d'erreur, statuts, enums. \\
\bottomrule
\end{tabularx}
\end{center}

\textbf{Chaîne complète d'une requête protégée :}
\begin{enumerate}
  \item Le frontend appelle un endpoint via \texttt{HttpClient}.
  \item L'intercepteur Angular ajoute le \texttt{Bearer token}.
  \item Spring Security intercepte la requête avant le contrôleur.
  \item Le JWT est vérifié ; l'identité utilisateur est extraite du token.
  \item Le contrôleur appelle le service du domaine concerné.
  \item Le service applique les règles métier et utilise un ou plusieurs repositories.
  \item Les données sont relues ou persistées en base.
  \item Le backend construit la réponse JSON finale et renvoie le code HTTP adapté
        (\texttt{200}, \texttt{201}, \texttt{204}, \texttt{400}, \texttt{401}, etc.).
\end{enumerate}

\textbf{Point important :} le frontend \textbf{n'envoie jamais l'identité utilisateur}
dans ses payloads métier. Le backend relie chaque opération au bon utilisateur à
partir du JWT validé. Cette décision simplifie le contrat, réduit le risque
d'usurpation et garantit que les données lues ou modifiées appartiennent bien
à la session authentifiée.

\subsection{Comment les éléments backend sont reliés entre eux}

Les routes visibles dans le frontend montrent une séparation nette par domaine.
Chaque domaine est très probablement raccordé à un ensemble cohérent de classes
backend reliées entre elles :

\begin{itemize}
  \item \textbf{Transactions} : un contrôleur reçoit les CRUD et recherches, un
        service applique les règles de création/mise à jour, puis un repository
        persiste les lignes de transaction.
  \item \textbf{Catégories} : les transactions référencent une catégorie ; le
        backend doit donc garantir qu'une catégorie utilisée appartient au même
        utilisateur et reste cohérente avec le type revenu/dépense.
  \item \textbf{Budgets} : les budgets sont liés aux catégories et à une période.
        Le backend agrège ou compare les transactions du mois avec les plafonds.
  \item \textbf{Objectifs} et \textbf{dettes} : chaque entité principale possède un
        historique paginé (\texttt{GoalContribution}, \texttt{DebtPayment}), ce qui
        implique un lien parent-enfants côté persistance et des services capables
        de recalculer les soldes après chaque mouvement.
  \item \textbf{Transactions planifiées} : elles produisent des exécutions
        (\texttt{PlannedTransaction}\allowbreak\texttt{Run}) et sont connectées au calendrier ainsi qu'au
        flux transactionnel réel lorsqu'une occurrence est lancée.
  \item \textbf{Calendrier} : il agrège les événements mensuels visibles dans l'UI ;
        il dépend donc soit de ses propres données, soit de projections calculées à
        partir des transactions planifiées.
  \item \textbf{Compte / sécurité} : \texttt{/account/logout-all} ne manipule pas les
        données métier ; il s'appuie sur l'identité extraite du token pour invalider
        les sessions côté serveur.
\end{itemize}

En pratique, la liaison entre ces éléments suit presque toujours le même axe :
\textbf{Controller} $\rightarrow$ \textbf{Service} $\rightarrow$ \textbf{Repository}
$\rightarrow$ \textbf{base de données}, puis la réponse remonte en sens inverse.
Le frontend ne connaît que le \textbf{contrat HTTP} ; tout le reste est encapsulé
dans le backend, ce qui permet de faire évoluer la logique métier sans casser
l'interface Angular tant que les DTO restent compatibles.

\subsection{Modèles partagés (\texttt{common.model.ts})}

\begin{lstlisting}
// Types union -- correspondent aux enums backend
export type TransactionType   = 'INCOME' | 'EXPENSE';
export type PlannedFrequency  = 'DAILY' | 'WEEKLY' | 'MONTHLY' | 'YEARLY';
export type GoalStatus        = 'ACTIVE' | 'COMPLETED' | 'ARCHIVED';
export type DebtStatus        = 'ACTIVE' | 'PAID' | 'CANCELLED';
export type PlannedRunStatus  = 'SUCCESS' | 'FAILED';

// Modèle d'erreur standardisé retourné par le backend
export interface ApiError {
  timestamp: string;
  status: number;
  error: string;
  message: string;
  path: string;
  details?: string[];
}

// Pagination Spring (page 0-indexée)
export interface Page<T> {
  content: T[];
  totalElements: number;
  totalPages: number;
  number: number;  // page courante (0-based)
  size: number;
}
\end{lstlisting}

\subsection{Validateur custom (\texttt{custom-validators.ts})}

Factory de validateur Angular utilisée sur les formulaires de changement de mot de passe :

\begin{lstlisting}
export function passwordMatchValidator(
  passwordField = 'password',
  confirmField  = 'confirmPassword',
): ValidatorFn {
  return (ctrl: AbstractControl): ValidationErrors | null => {
    const pw  = ctrl.get(passwordField)?.value;
    const cpw = ctrl.get(confirmField)?.value;
    return pw && cpw && pw !== cpw ? { passwordMismatch: true } : null;
  };
}
\end{lstlisting}

C'est une \textbf{factory} (fonction qui retourne un \texttt{ValidatorFn}) afin de
configurer les noms des champs à comparer selon le formulaire :
\begin{itemize}
  \item \texttt{signup-page} et \texttt{reset-password-page} : champs \texttt{password}
        + \texttt{confirmPassword} (valeurs par défaut)
  \item \texttt{profile-page} : champs \texttt{newPassword} + \texttt{confirmPassword}
\end{itemize}

% ════════════════════════════════════════════════════════════════════════════
\section{Couche Shared -- Composants réutilisables}

\subsection{PaginationComponent}

Composant générique de pagination qui consomme le format \texttt{Page<T>} de Spring :

\begin{lstlisting}
@Component({ selector: 'app-pagination', standalone: true, ... })
export class PaginationComponent {
  @Input() page         = 0;          // 0-based (Spring)
  @Input() totalPages   = 0;
  @Input() totalElements = 0;
  @Input() size         = 20;
  @Output() pageChange  = new EventEmitter<number>();
}
\end{lstlisting}

\textbf{Comportement :}
\begin{itemize}
  \item Affiche \og N--M sur TOTAL \fg{} (converti en 1-based pour l'utilisateur).
  \item Génère une liste intelligente de boutons de page : toujours les 2 premières
        et les 2 dernières pages, plus une fenêtre de 2 pages autour de la page courante.
  \item Insère des points de suspension (\texttt{...}) lorsqu'il y a un saut.
  \item Boutons Précédent / Suivant désactivés aux extrémités.
  \item Émet l'événement \texttt{pageChange} sur sélection d'une page.
\end{itemize}

% ════════════════════════════════════════════════════════════════════════════
\section{Couche Layout}

\subsection{MainLayout}

Conteneur de toutes les pages privées (\texttt{/app/**}) :

\begin{lstlisting}
@Component({
  selector: 'app-main-layout',
  standalone: true,
  imports: [RouterOutlet, AppSidebar],
})
export class MainLayout {}
\end{lstlisting}

Le template positionne la sidebar à gauche (largeur fixe 256 px) et le
contenu principal à droite :

\begin{lstlisting}[language=html, numbers=none]
<div class="flex min-h-screen bg-[#f6f6f8]">
  <app-sidebar />
  <div class="ml-64 flex-1">
    <router-outlet />
  </div>
</div>
\end{lstlisting}

\subsection{AppSidebar (\texttt{features/sidebar/})}

Composant de navigation latérale chargé par \texttt{MainLayout}.

\textbf{Structure de la navigation (sections)} :
\begin{itemize}
  \item \textit{Vue d'ensemble} : Tableau de bord, Transactions, Planifiées
  \item \textit{Gestion} : Catégories, Budgets, Objectifs, Dettes
  \item \textit{Outils} : Calendrier, Import / Export
  \item \textit{Compte} : Paramètres (profile)
\end{itemize}

\textbf{Fonctionnalités :}
\begin{itemize}
  \item Affiche le nom complet et les initiales de l'utilisateur (lues depuis les
        métadonnées Supabase).
  \item Menu déroulant utilisateur fermé par \texttt{HostListener('document:click')}.
  \item Déconnexion locale via SweetAlert2 $\to$ \texttt{authService.signOut()} $\to$
        redirection vers \texttt{/}.
\end{itemize}

% ════════════════════════════════════════════════════════════════════════════
\section{Couche Features -- Services et modèles par domaine}

Chaque domaine métier suit la même organisation :
\begin{lstlisting}[language=bash, numbers=none, frame=none, backgroundcolor=\color{codebg}, basicstyle=\ttfamily\small\color{codefg}]
features/<domaine>/
  models/<domaine>.model.ts    # Interfaces TypeScript
  services/<domaine>.service.ts # Appels HTTP
  pages/<domaine>-page/
    <domaine>-page.ts          # Logique du composant
    <domaine>-page.html        # Template
    <domaine>-page.css/scss
\end{lstlisting}

\subsection{Transactions}

\subsubsection{Modèles}

\begin{lstlisting}
export interface Transaction {
  id: number;         label: string;      amount: number;
  type: TransactionType;  categoryId: number; date: string; // YYYY-MM-DD
  note?: string | null;   source?: string | null;
  createdAt: string;      updatedAt?: string | null;
}

export interface CreateTransactionRequest {
  amount: number;   type: TransactionType;  categoryId: number;
  label: string;    date: string;           note?: string | null;
}

export interface TransactionSearchParams {
  type?: TransactionType;  categoryId?: number;
  startDate?: string;      endDate?: string;    label?: string;
}
\end{lstlisting}

\subsubsection{Service}

\begin{lstlisting}
@Injectable({ providedIn: 'root' })
export class TransactionsService {
  private base = `${environment.apiBaseUrl}/transactions`;

  getAll(page = 0, size = 200): Observable<Page<Transaction>>
  search(params: TransactionSearchParams, page = 0, size = 200):
      Observable<Page<Transaction>>
  create(payload: CreateTransactionRequest): Observable<Transaction>
  update(id: number, payload: UpdateTransactionRequest): Observable<Transaction>
  delete(id: number): Observable<void>
}
\end{lstlisting}

\subsubsection{Page}

Charge \textbf{toutes} les transactions en mémoire (\texttt{size=200}),
puis applique le filtrage et la pagination \textbf{côté client}. Ce choix simplifie
l'implémentation pour un projet académique à faible volumétrie.

\subsection{Catégories}

\begin{lstlisting}
export interface Category {
  id: number;   name: string;
  icon?: string | null;   color?: string | null;
  isDefault: boolean;     createdAt: string;
}
\end{lstlisting}

La page distingue deux onglets (\textit{Dépenses} / \textit{Revenus}) et propose :
\begin{itemize}
  \item Une banque de 60+ icônes Material Symbols avec recherche textuelle.
  \item 12 couleurs prédéfinies pour l'icône.
  \item La suppression est bloquée sur les catégories \texttt{isDefault = true}.
\end{itemize}

\subsection{Budgets}

\begin{lstlisting}
export interface Budget {
  id: number;           categoryId: number;
  monthlyLimit: number; alertThreshold: number; // 0-100
  createdAt: string;    updatedAt?: string | null;
}
\end{lstlisting}

La page enrichit chaque budget avec le nom et la couleur de la catégorie associée.
Le seuil d'alerte (\texttt{alertThreshold}) est un pourcentage : si le montant
dépensé atteint ce seuil, le budget s'affiche en orange ou rouge.

\subsection{Objectifs (\texttt{goals})}

\begin{lstlisting}
export interface Goal {
  id: number;           name: string;
  targetAmount: number; currentAmount: number;
  deadline?: string | null;  status: GoalStatus;
}

export interface GoalContribution {
  id: number;           goalId: number;
  amount: number;       contributionDate: string;
  note?: string | null; transactionId?: number | null;
}

export interface ContributeGoalRequest {
  amount: number;  contributionDate: string;  note?: string | null;
}
\end{lstlisting}

\textbf{Endpoints spéciaux :}
\begin{itemize}
  \item \texttt{PATCH /v1/goals/\{id\}/contribute} -- ajoute une contribution.
  \item \texttt{GET /v1/goals/\{id\}/contributions} -- historique paginé.
\end{itemize}

La page affiche des cercles de progression via un \texttt{conic-gradient} CSS calculé
dynamiquement. Les contributions sont chargées \textbf{à la demande} (\textit{lazy})
par objectif.

\subsection{Dettes}

\begin{lstlisting}
export interface Debt {
  id: number;              creditorName: string;
  totalAmount: number;     remainingAmount: number;
  monthlyPayment: number;  interestRate?: number | null;
  startDate: string;       dueDate?: string | null;
  status: DebtStatus;
}

export interface DebtSummary {
  activeDebtsCount: number;   paidDebtsCount: number;
  totalRemainingAmount: number; totalMonthlyPayment: number;
}
\end{lstlisting}

\textbf{Endpoints spéciaux :}
\begin{itemize}
  \item \texttt{PATCH /v1/debts/\{id\}/payment} -- enregistre un paiement.
  \item \texttt{GET /v1/debts/\{id\}/payments} -- historique paginé.
  \item \texttt{GET /v1/debts/summary} -- agrégats globaux.
\end{itemize}

\subsection{Transactions planifiées}

\begin{lstlisting}
export interface PlannedTransaction {
  id: number;          label: string;     amount: number;
  type: TransactionType; categoryId: number;
  frequency: PlannedFrequency;  nextDueDate: string;
  autoCreate: boolean; active: boolean;
}
\end{lstlisting}

\textbf{Endpoints spéciaux :}
\begin{itemize}
  \item \texttt{PATCH /v1/planned-transactions/\{id\}/toggle} -- active/désactive.
  \item \texttt{POST /v1/planned-transactions/\{id\}/run-now} -- exécution manuelle.
  \item \texttt{GET /v1/planned-transactions/\{id\}/runs} -- historique d'exécutions.
\end{itemize}

La page applique plusieurs règles UI spécifiques :
\begin{itemize}
  \item La carte \og Prochaine échéance \fg{} trie les planifications filtrées
        par \texttt{nextDueDate} croissant et affiche la plus proche échéance,
        au lieu de reprendre le premier élément brut renvoyé par l'API.
  \item Le formulaire de création/édition force les champs
        \texttt{amount}, \texttt{type}, \texttt{frequency},
        \texttt{nextDueDate} et \texttt{autoCreate} avec des contrôles
        \texttt{nonNullable}, pour ne jamais envoyer \texttt{null} sur ces
        valeurs au backend.
  \item Après une exécution manuelle réussie (\texttt{run-now}), la page recharge
        la liste complète afin de récupérer le \texttt{nextDueDate} recalculé
        côté backend et de rafraîchir immédiatement la carte ainsi que le tableau.
\end{itemize}

\subsection{Calendrier}

\begin{lstlisting}
export interface CalendarEvent {
  id: number;   title: string;   eventDate: string; // YYYY-MM-DD
  amount?: number | null;   color?: string | null;
  notes?: string | null;
}
\end{lstlisting}

La page construit une grille de 42 jours (6 semaines × 7). Les transactions
planifiées actives sont \textbf{projetées} dans la grille selon leur fréquence :
\begin{itemize}
  \item \texttt{DAILY} -- apparaît chaque jour.
  \item \texttt{WEEKLY} -- même jour de semaine toutes les 7 jours.
  \item \texttt{MONTHLY} -- même date chaque mois.
  \item \texttt{YEARLY} -- même date et même mois.
\end{itemize}

\subsection{Tableau de bord}

\begin{lstlisting}
export interface MonthlySummary {
  month: string; totalIncome: number;
  totalExpenses: number; netBalance: number;
  transactionCount: number;
}

export interface SpendingByCategory {
  categoryId: number; categoryName: string;
  amount: number; percentage: number;
}

export interface GoalProgress {
  id: number; name: string;
  targetAmount: number; currentAmount: number;
  percentage: number; status: GoalStatus;
}

export interface BudgetAlert {
  categoryName: string; monthlyLimit?: number;
  spentAmount?: number; percentage?: number;
  overBudget?: boolean;
}
\end{lstlisting}

La page lance un \texttt{forkJoin} qui déclenche \textbf{7 requêtes HTTP en parallèle}
au chargement :

\begin{lstlisting}
forkJoin({
  summary:   dashboardService.getSummary(month),
  spending:  dashboardService.getSpendingByCategory(month),
  planned:   plannedService.getAll(),
  goals:     dashboardService.getGoalsProgress(),
  alerts:    dashboardService.getBudgetAlerts(month),
  recent:    transactionsService.getAll(0, 5),
  monthlyTx: transactionsService.search({ startDate, endDate }, 0, 500),
}).subscribe({ next: (data) => { /* populate widgets */ } });
\end{lstlisting}

Chaque requête est protégée par \texttt{catchError(() => of(valeurParDefaut))} afin
qu'une panne partielle ne bloque pas l'ensemble du tableau de bord.

\textbf{Le seul graphique chart.js} de l'application (ligne chronologique) est ici :
\begin{lstlisting}
lineData: ChartData<'line'> = { labels: [], datasets: [{
  data: [], borderColor: '#0f766e', fill: true, tension: 0.4
}]};
// Dans le template :
// <canvas baseChart [data]="lineData" [options]="lineOptions" type="line">
\end{lstlisting}

Le graphique "donut" de répartition des dépenses est du \textbf{CSS pur} :
\begin{lstlisting}
get spendingBackground(): string {
  // Construit un conic-gradient avec les proportions des catégories
  return `conic-gradient(${stops.join(',')})`;
}
\end{lstlisting}

\subsection{Import / Export CSV}

\textbf{Export :} récupère toutes les transactions via \texttt{getAll(0,10000)},
génère un \texttt{Blob} CSV avec BOM UTF-8 (pour compatibilité Excel) et déclenche
un téléchargement via un lien \texttt{<a>} créé dynamiquement.

\textbf{Import :} lit le fichier CSV via \texttt{FileReader}, valide chaque ligne,
puis appelle \texttt{create()} pour chaque ligne valide.
Les erreurs de parsing et les erreurs HTTP sont collectées et affichées dans un
rapport détaillé.

Format CSV attendu :
\begin{lstlisting}[language=bash, numbers=none]
date,label,amount,type,categoryId,note
2025-03-01,Courses,45.50,EXPENSE,3,Supermarché
2025-03-02,Salaire,2500,INCOME,1,
\end{lstlisting}

\subsection{Profil}

Le service \texttt{ProfileService} appelle directement le \textbf{SDK Supabase}
(pas le backend Spring) pour mettre à jour le nom et le mot de passe :

\begin{lstlisting}
updateProfile(req: ProfileUpdateRequest): Promise<User> {
  return supabase.auth.updateUser({
    data: { full_name: req.fullName }
  });
}

updatePassword(req: PasswordUpdateRequest): Promise<User> {
  return supabase.auth.updateUser({ password: req.newPassword });
}
\end{lstlisting}

La page \texttt{ProfilePage} inclut également la déconnexion globale via
\texttt{AccountService} (situé dans \texttt{core/services/}), qui appelle
\texttt{POST /v1/account/logout-all} sur le backend Spring pour invalider
toutes les sessions serveur.

% ════════════════════════════════════════════════════════════════════════════
\section{Flux de données -- exemple complet}

\begin{mdframed}[style=infobox]
\textbf{Exemple : création d'une transaction}
\end{mdframed}

\begin{enumerate}
  \item L'utilisateur remplit le formulaire Reactive Forms et clique sur \og Créer \fg.
  \item Le composant vérifie \texttt{form.invalid} : si invalide, tous les champs
        sont marqués \texttt{touched} pour afficher les erreurs ; on retourne.
  \item \texttt{mapFormToRequest()} construit le payload :
        un \texttt{CreateTransactionRequest} typé.
  \item \texttt{transactionsService.create(payload)} est appelé.
        Il retourne un \texttt{Observable<}\allowbreak\texttt{Transaction>}.
  \item La requête \texttt{POST /v1/transactions} est construite
        par \texttt{HttpClient}.
  \item \texttt{authInterceptor} lit le JWT depuis le cache et ajoute
        \texttt{Authorization: Bearer <token>}.
  \item Le filtre de sécurité Spring valide le JWT et rattache la requête à
        l'utilisateur authentifié.
  \item Le contrôleur REST reçoit le JSON, le service métier vérifie les règles
        de création, puis le repository persiste la transaction en base.
  \item Le backend retourne \texttt{201 Created} avec le \texttt{Transaction} JSON
        sérialisé pour le frontend.
  \item \texttt{HttpClient} désérialise en \texttt{Transaction}.
  \item Callback \texttt{next} : fermeture du formulaire, rechargement de la liste,
        toast SweetAlert2 succès (\texttt{timer: 1500}).
  \item Callback \texttt{error} : toast SweetAlert2 erreur avec message explicite.
\end{enumerate}

\begin{mdframed}[style=infobox]
\textbf{Exemple : exécution manuelle d'une transaction planifiée}
\end{mdframed}

\begin{enumerate}
  \item L'utilisateur clique sur l'action \texttt{run-now} depuis une ligne de la
        page \texttt{PlannedPage}.
  \item Une confirmation SweetAlert2 sécurise l'action avant l'appel HTTP.
  \item \texttt{plannedService.runNow(id)} envoie
        \texttt{POST /v1/planned-transactions/\{id\}/}\allowbreak\texttt{run-now}.
  \item Le backend crée la transaction réelle, enregistre une ligne
        \texttt{PlannedTransactionRun}, puis avance \texttt{nextDueDate} selon
        la fréquence de la planification.
  \item En cas de succès, le frontend recharge les planifications et, si le détail
        est ouvert, recharge aussi l'historique des exécutions.
  \item La colonne \texttt{nextDueDate} du tableau et la carte \og Prochaine
        échéance \fg{} affichent alors immédiatement la nouvelle date renvoyée
        par le backend.
\end{enumerate}

% ════════════════════════════════════════════════════════════════════════════
\section{Flux d'authentification}

\subsection{Connexion}

\begin{enumerate}
  \item L'utilisateur soumet le formulaire \texttt{LoginPage}.
  \item \texttt{authService.signIn(email, password)} appelle le SDK Supabase.
  \item Supabase retourne un \texttt{Session} (contenant le JWT).
  \item \texttt{BehaviorSubject} est mis à jour via l'événement \texttt{onAuthStateChange}.
  \item Le router navigue vers \texttt{/app/dashboard}.
  \item Désormais, chaque requête HTTP inclura automatiquement le token via l'intercepteur.
\end{enumerate}

\subsection{Accès à une route protégée}

\begin{enumerate}
  \item L'utilisateur accède à \texttt{/app/transactions}.
  \item \texttt{authGuard} s'exécute : \texttt{await auth.waitForInit()}.
  \item Si session absente $\to$ redirection vers \texttt{/login}.
  \item Si session présente $\to$ la route est activée, \texttt{MainLayout} est rendu.
\end{enumerate}

\subsection{Déconnexion}

\begin{enumerate}
  \item L'utilisateur clique sur \og Se déconnecter \fg{} dans la sidebar.
  \item SweetAlert2 demande confirmation.
  \item \texttt{authService.signOut()} vide la session Supabase locale.
  \item \texttt{BehaviorSubject} émet \texttt{null}.
  \item Le router redirige vers \texttt{/}.
  \item L'intercepteur ne trouvant plus de token, les prochaines requêtes seront anonymes.
\end{enumerate}

% ════════════════════════════════════════════════════════════════════════════
\section{Formulaires Reactive Forms}

Tous les formulaires métier utilisent \textbf{Angular Reactive Forms}.
Le pattern appliqué dans chaque composant est :

\begin{lstlisting}
// Déclaration du groupe
form = this.fb.group({
  label:      this.fb.nonNullable.control('', [Validators.required, Validators.maxLength(80)]),
  amount:     this.fb.nonNullable.control(0, [Validators.required, Validators.min(0.01)]),
  type:       this.fb.nonNullable.control('EXPENSE' as TransactionType, [Validators.required]),
  categoryId: [null as number | null, [Validators.required]],
  frequency:  this.fb.nonNullable.control('MONTHLY' as PlannedFrequency, [Validators.required]),
  nextDueDate:this.fb.nonNullable.control(dayjs().format('YYYY-MM-DD'), [Validators.required]),
  autoCreate: this.fb.nonNullable.control(false),
});

// Accesseurs typés (évite les .get('label') répétés dans le template)
get label()      { return this.form.controls.label; }
get amount()     { return this.form.controls.amount; }

// Soumission
submit(): void {
  if (this.form.invalid) { this.form.markAllAsTouched(); return; }
  const payload = this.mapFormToRequest();
  this.service.create(payload).subscribe({ ... });
}
\end{lstlisting}

Cette approche réduit les erreurs de sérialisation côté backend : les champs
primitifs ou enum ne partent plus avec \texttt{null} dans le JSON, ce qui est
particulièrement important pour des propriétés comme \texttt{boolean autoCreate}.

Le template exploite les états Angular pour l'UX :

\begin{lstlisting}[language=html, numbers=none]
<input formControlName="label"
       [class.border-red-500]="label.invalid && label.touched"
       required maxlength="80" />
@if (label.invalid && label.touched) {
  <p class="text-red-500 text-sm mt-1">
    @if (label.hasError('required'))  { Le libellé est requis. }
    @if (label.hasError('maxlength')) { 80 caractères maximum. }
  </p>
}
<button [disabled]="form.invalid || saving" type="submit">
  Créer
</button>
\end{lstlisting}

% ════════════════════════════════════════════════════════════════════════════
\section{Référence des endpoints HTTP}

Base URL : \texttt{http://localhost:8080/v1}

\begin{longtable}{>{\small}L{1.9cm} >{\small\ttfamily}L{7.4cm} >{\small\ttfamily}L{4.5cm}}
\toprule
\textbf{Méthode} & \textbf{Endpoint} & \textbf{Retour} \\
\midrule
GET & /transactions & Page\textless Transaction\textgreater \\
GET & /transactions/search & Page\textless Transaction\textgreater \\
POST & /transactions & Transaction \\
PUT & /transactions/\{id\} & Transaction \\
DELETE & /transactions/\{id\} & void \\
\midrule
GET & /categories & Category[] \\
POST & /categories & Category \\
PUT & /categories/\{id\} & Category \\
DELETE & /categories/\{id\} & void \\
\midrule
GET & /budgets & Budget[] \\
POST & /budgets & Budget \\
PUT & /budgets/\{id\} & Budget \\
DELETE & /budgets/\{id\} & void \\
\midrule
GET & /goals & Goal[] \\
POST & /goals & Goal \\
PUT & /goals/\{id\} & Goal \\
DELETE & /goals/\{id\} & void \\
PATCH & /goals/\{id\}/contribute & Goal \\
GET & /goals/\{id\}/contributions & Page\textless GoalContribution\textgreater \\
\midrule
GET & /debts & Debt[] \\
POST & /debts & Debt \\
PUT & /debts/\{id\} & Debt \\
DELETE & /debts/\{id\} & void \\
PATCH & /debts/\{id\}/payment & Debt \\
GET & /debts/\{id\}/payments & Page\textless DebtPayment\textgreater \\
GET & /debts/summary & DebtSummary \\
\midrule
GET & /planned-transactions & PlannedTransaction[] \\
POST & /planned-transactions & PlannedTransaction \\
PUT & /planned-transactions/\{id\} & PlannedTransaction \\
DELETE & /planned-transactions/\{id\} & void \\
PATCH & /planned-transactions/\{id\}/toggle & PlannedTransaction \\
POST & /planned-transactions/\{id\}/run-now & PlannedTransactionRun \\
GET & /planned-transactions/\{id\}/runs & Page\textless{}\allowbreak PlannedTransactionRun\textgreater{} \\
\midrule
GET & /calendar/events?month=YYYY-MM & CalendarEvent[] \\
POST & /calendar/events & CalendarEvent \\
PUT & /calendar/events/\{id\} & CalendarEvent \\
DELETE & /calendar/events/\{id\} & void \\
\midrule
GET & /dashboard/summary?month=YYYY-MM & MonthlySummary \\
GET & /dashboard/spending-by-category & SpendingByCategory[] \\
GET & /dashboard/goals-progress & GoalProgress[] \\
GET & /dashboard/budget-alerts & BudgetAlert[] \\
\midrule
POST & /account/logout-all & void \\
\bottomrule
\end{longtable}

% ════════════════════════════════════════════════════════════════════════════
\section{Décisions d'architecture}

\begin{center}
\begin{tabularx}{\linewidth}{L{5cm} X}
\toprule
\textbf{Décision} & \textbf{Justification} \\
\midrule
Standalone components & Plus de NgModule ; chaque composant déclare ses propres
                        imports. Plus lisible, plus modulaire, conforme aux bonnes
                        pratiques Angular 17+. \\
\midrule
Lazy loading sur toutes les routes & Le bundle initial est minimal. Le code d'une
                                     page n'est téléchargé que lorsque l'utilisateur
                                     y accède. \\
\midrule
BehaviorSubject pour la session & Lecture synchrone du token dans l'intercepteur.
                                   Pas besoin de souscrire à un Observable à chaque
                                   requête. \\
\midrule
Un service par domaine & Séparation stricte des responsabilités. Chaque service
                          connaît uniquement ses propres endpoints. \\
\midrule
Filtrage côté client (transactions) & Simplification pour une volumétrie faible.
                                       Les endpoints de recherche backend existent
                                       pour un usage en production. \\
\midrule
Pas d'intercepteur d'erreur centralisé & Chaque page gère ses erreurs avec le contexte
                                          métier approprié (message adapté à l'action). \\
\midrule
Reactive Forms uniquement & Validation complexe (cross-fields, conditions) impossible
                            proprement en template-driven forms. \\
\midrule
SweetAlert2 & UX supérieure à \texttt{window.alert/confirm}. Couleurs cohérentes
              avec la charte graphique. \\
\midrule
Pas d'envoi de \texttt{userId} & Le backend identifie l'utilisateur depuis le JWT.
                                  Aucun identifiant ne transite dans le corps des requêtes. \\
\bottomrule
\end{tabularx}
\end{center}

% ════════════════════════════════════════════════════════════════════════════
\section{Sécurité}

\begin{itemize}
  \item La clé Supabase \texttt{anon} est la seule clé présente côté frontend.
        Elle est publique par définition et ne donne accès qu'aux ressources
        que les règles RLS (Row Level Security) Supabase autorisent.
  \item La clé \texttt{service\_role} (clé admin Supabase) n'est \textbf{jamais} exposée
        côté frontend.
  \item Le JWT n'est jamais loggé ni stocké en dehors du cache interne du SDK Supabase
        (localStorage ou cookies selon la configuration Supabase).
  \item Le backend Spring valide le JWT sur chaque requête protégée et extrait
        l'identité de l'utilisateur depuis le token, jamais depuis le corps de la requête.
  \item \texttt{logout-all} appelle \texttt{POST /v1/account/logout-all} sans
        paramètre \texttt{userId} : le backend déduit l'utilisateur depuis le token.
\end{itemize}

% ════════════════════════════════════════════════════════════════════════════
\section{Environnements}

\begin{lstlisting}
// environment.ts (développement)
export const environment = {
  production:      false,
  apiBaseUrl:      'http://localhost:8080/v1',
  supabaseUrl:     'https://YOUR_PROJECT.supabase.co',
  supabaseAnonKey: 'YOUR_SUPABASE_ANON_KEY',
};
\end{lstlisting}

Toutes les URL backend sont construites relativement à \texttt{environment.apiBaseUrl}.
Il n'existe \textbf{qu'un seul préfixe} dans l'application : \texttt{/v1}.

% ════════════════════════════════════════════════════════════════════════════
\section{Conclusion}

Le frontend Finance Tracker est une SPA Angular bien structurée, strictement typée
et sécurisée. Les principes de conception retenus -- séparation claire des couches,
lazy loading, Reactive Forms, singleton services, token JWT via intercepteur --
en font une base maintenable et extensible.

Chaque domaine métier est autonome : ses modèles, son service et ses composants
sont colocalisés dans \texttt{features/<domaine>/}. Les parties transverses
(auth, guards, intercepteurs, types communs) sont isolées dans \texttt{core/}.
Les composants réutilisables vivent dans \texttt{shared/}.

Cette architecture est conçue pour rester lisible et explicable à mesure que
le projet évolue, sans sur-ingénierie (pas de state management externe, pas de
multi-comptes, pas de devise multiple).

\end{document}
